\lesson{3}{Sep 10 2021 Fri (06:53:22)}{Energy and Temperature}{Unit 1}

\subsubsection*{Different Forms of Energy}

Particles in motion can cause change due to their motion. For example:

\begin{definition}[Energy]
	\bf{Energy} is the ability to cause change or the ability to do work.
	\bf{Energy} is measured in the \bf{SI Unit Joules}:

    \begin{align}
    	\rm{Joule} = \frac{kg \times m^2}{s^^2}
    \end{align}
\end{definition}

\begin{definition}[Work]
	\bf{Work} is the process of causing matter to move against an opposing
	force. When matter moves, it's position is changing. For that to happen, it
	requires \bf{Energy}.
\end{definition}

All forms of energy can be broadly classified as either:

\begin{itemize}
	\item \bf{Potential Energy}.
	\item \bf{Kinetic Energy}.
\end{itemize}

\begin{definition}[Potential Energy]
	\bf{Potential Energy} is the energy an object has because of it's current
	position or composition. For example, you have someone standing on a plane
	$25,000 ft$ above ground before he sky dives which gives in \bf{Potential Energy}.

	\bf{Potential Energy} is also known as \bf{Stored Energy} because
	it's energy that is there, but hasn't been used. You can think about it as,
	winding up a toy car, then once you let it go, all of the energy comes out as
	\bf{Kinetic Energy}.
\end{definition}

\begin{definition}[Kinetic Energy]
	\bf{Kinetic Energy} is energy an object
	has due to its motion. Like I mentioned before, when you wind up a toy car
	and then let it go, that's an example of \bf{Kinetic Energy}.
\end{definition}

\begin{definition}[Mechanical Energy]
	\bf{Mechanical Energy} is the movement of objects from one place to
	another

	EXAMPLES:
	\begin{itemize}
		\item Turning wheels.
		\item Wind blowing.
	\end{itemize}
\end{definition}

\begin{definition}[Stored Mechanical Energy]
	\bf{Stored Mechanical Energy} is the energy stored in objects by
	application of force.

	EXAMPLES:
	\begin{itemize}
		\item Compressed spring.
		\item Pulled back string of a bow before shooting the arrow.
	\end{itemize}
\end{definition}

\begin{definition}[Electrical Energy]
	\bf{Electrical Energy} is the movement of electrical charge.

	EXAMPLES:
	\begin{itemize}
		\item Electricity running though wires.
		\item Lighting.
	\end{itemize}
\end{definition}

\begin{definition}[Electrical Potential Energy]
	\bf{Electrical Potential Energy} is the stored energy of electric charges
	due to their position and interaction around other charges.

	EXAMPLES:
	\begin{itemize}
		\item Non-moving charges in circuits.
	\end{itemize}
\end{definition}

\begin{definition}[Radiant Energy]
	\bf{Radiant Energy} is electromagnetic energy that travels in waves.

	EXAMPLES:
	\begin{itemize}
		\item X-Rays.
		\item Microwaves.
		\item Visible light.
		\item Radio waves.
	\end{itemize}
\end{definition}

\begin{definition}[Chemical Energy]
	\bf{Chemical Energy} is stored in the bonds of atoms and molecules.

	EXAMPLES:
	\begin{itemize}
		\item Biomass.
		\item Natural gas.
		\item Propane.
	\end{itemize}
\end{definition}

\begin{definition}[Sound Energy]
	\bf{Sound Energy} is the movement of energy through substances in
	longitudinal waves.  Sound is produced when a force causes an object to move
	or vibrate, which is when energy is transferred through the object in a wave.
\end{definition}

\begin{definition}[Nuclear Energy]
	\bf{Nuclear Energy} is energy stored in the nucleus of an atom. This is
	the energy that holds atoms together and can be released in processes called
	fusion and fission.
\end{definition}

\begin{definition}[Thermal Energy]
	\bf{Thermal Energy} also known as \bf{heat}, \bf{Thermal Energy}
	is the internal energy in a substance caused by vibration of atoms and
	molecules. \bf{Geothermal Energy} is an example of \bf{Thermal
	Energy}. Heat is also a byproduct of a lot of energy conversions.
\end{definition}

\begin{definition}[Gravitational Energy]
	\bf{Gravitational Energy} is the \bf{Potential Energy} of position.
	The rock resting at the top of a hill has \bf{Gravitational Potential
	Energy}.
\end{definition}

Although there are many forms of energy, the \bf{Law of Conservation of Energy} states that
\it{\bf{"Energy can be converted from one form to another, but it's not created or destroyed in ordinary physical and chemical processes."}}

\subsubsection*{Energy being conserved during Physical and Chemical Change}

Energy rarely stays in its current form. Instead, it is converted to different
forms through \bf{physical} and \bf{chemical} processes. As it changes
from one form to another within these processes, more energy is not created, and
none of the energy is destroyed; it simply changes form.

\subsubsection*{Heat Vs. Temperature}

To understand the relationship between heat and the motion of particles, let's
look at how thermal energy affects the interaction between and the motion of
sugar and coffee particles.

For the sugar crystals to dissolve, the crystals need to break apart into
smaller molecules. When we stir, we help break down the solid particles by
causing them to move and change faster. When heat transfers to the molecules in
the coffee, they vibrate faster, increasing their kinetic energy. This increases
the number of collisions between the particles to break down sugar molecules.

The average kinetic energy of molecules is what determines the temperature of a
substance. So, hotter objects have faster-moving particles and higher
temperatures. Temperature measures the averages because moving particles can
absorb and release heat with each collision as heat is converted to kinetic
energy and back.

\newpage
